% Part: first-order-logic
% Chapter: syntax-and-semantics
% Section: unique-readability

\documentclass[../../../include/open-logic-section]{subfiles}

\begin{document}

\olfileid{fol}{syn}{unq}

\olsection{એકમાત્ર વાચનીયતા}

\begin{explain}
આપણે જે રીતે !!{formula}s વ્યાખ્યાયિત કર્યાં છે તેનાથી દરેક
!!{formula}નું \emph{એકમાત્ર વાંચન} હોવાની ખાતરી મળે છે; એટલે કે,
!!{formula}s માટેના આપણા રચના-નિયમો અનુસાર તેને રચવાની મૂળભૂત રીતે
માત્ર એક જ રીત અને તેનો ``અર્થઘટન'' કરવાની માત્ર એક જ રીત હોય છે. જો
એમ ન હોત તો આપણી પાસે સંદિગ્ધ !!{formula}s હોત, એટલે કે એકથી વધુ
વાંચન અથવા અર્થઘટન ધરાવતાં !!{formula}s---અને આપણે દેખીતી રીતે એ
ટાળવા માગીએ છીએ. વધુ મહત્વની વાત એ છે કે આ ગુણધર્મ વિના આપણે આગળ
આપવાની મોટા ભાગની વ્યાખ્યાઓ અને સાબિતીઓ કામ નહિ કરે.

આ વાત કદાચ ત્યારે સૌથી સ્પષ્ટ થાય જ્યારે જોઈએ કે એકમાત્ર વાચનીયતાની
ખાતરી ન આપતા !!{formula}sના ખરાબ રચના-નિયમો આપ્યા હોત તો શું થાત.
દાખલા તરીકે, સંયોજકો માટેના રચના-નિયમોમાં કૌંસ મૂકવાનું ભૂલી ગયા હોત
અને આવું માન્ય રાખ્યું હોત:
\begin{quote}
જો $!A$ અને $!B$ !!{formula}s હોય, તો $!A \lif !B$ પણ સૂત્ર છે.
\end{quote}
આણ્વિક સૂત્ર~$!D$થી શરૂ કરતાં આ નિયમ આપણને $!D\lif !D$ રચવા દે.
તેમાંથી, $!D$ સાથે, $!D\lif !D\lif !D$ મળે. પરંતુ આ કરવાની બે રીતો છે:
\begin{enumerate}
\item આપણે $!D$ને $!A$ અને $!D\lif !D$ને $!B$ લઈએ.
\item આપણે $!A$ને $!D\lif !D$ અને $!B$ને~$!D$ લઈએ.
\end{enumerate}
એને અનુરૂપ, !!{formula}~$!D\lif !D\lif !D$ને ``વાંચવાની'' પણ બે રીતો
છે. તે $!B\lif !C$ સ્વરૂપનું છે, જ્યાં $!B$ એ $!D$ અને $!C$ એ
$!D\lif !D$ છે; પરંતુ તે \emph{સાથે સાથે} $!B\lif !C$ સ્વરૂપનું એવું
પણ છે જેમાં $!B$ એ $!D\lif !D$ અને $!C$ એ~$!D$ છે.

આવું બને તો આપણી વ્યાખ્યાઓ હંમેશાં કામ કરશે નહિ. દાખલા તરીકે,
સૂત્રનો !!{main operator} વ્યાખ્યાયિત કરતાં આપણે કહીએ છીએ:
$!B\lif !C$ સ્વરૂપના સૂત્રમાં દર્શાવેલું~$\lif$નું આવર્તન
!!{main operator} છે. પરંતુ $!D\lif !D\lif !D$ સૂત્રને ઉપર જણાવેલી બે
જુદી રીતોથી $!B\lif !C$ સાથે મેળવો, તો એક કિસ્સામાં $\lif$નું પહેલું
આવર્તન !!{main operator} બને અને બીજા કિસ્સામાં બીજું. જ્યારે આપણી
મનશા એવી છે કે !!{main operator}, !!{formula}નું \emph{વિધેય}
હોય; એટલે કે, દરેક !!{formula}માં !!{main
  operator}નું બરાબર એક
આવર્તન હોવું જોઈએ.
\end{explain}

\begin{lem}
!!a{formula}~$!A$માં ડાબા અને જમણા કૌંસોની સંખ્યા સમાન હોય છે.
\end{lem}

\begin{proof}
$!A$ જે રીતે રચાય છે તેના પર અનુમાન વડે આપણે આ સાબિત કરીએ છીએ.
એ માટે બે વસ્તુઓ જોઈએ: (ક) પહેલાં બતાવવું પડે કે બધાં આણ્વિક સૂત્રોમાં
આ ગુણધર્મ છે (અનુમાનનો આધાર). (ખ) પછી બતાવવું પડે કે આપેલાં
સૂત્રોમાંથી નવાં સૂત્રો રચીએ ત્યારે, જૂનાં સૂત્રોમાં ગુણધર્મ હોય તો
નવાંમાં પણ હોય છે.

$l(!A)$ને~$!A$માંના ડાબા કૌંસોની સંખ્યા અને $r(!A)$ને જમણા
કૌંસોની સંખ્યા કહો; એવી જ રીતે $l(t)$ અને $r(t)$ને પદ~$t$માંના
અનુક્રમે ડાબા અને જમણા કૌંસોની સંખ્યા કહો.

\begin{prob}
    કોઈ પણ પદ~$t$ માટે $l(t)=r(t)$ સાબિત કરો.
\end{prob}

\begin{enumerate}
\tagitem{prvFalse}{\indcase{!A}{\lfalse}{$\indfrm$માં $0$ ડાબા અને
    $0$ જમણા કૌંસ છે.}}{}

\tagitem{prvTrue}{\indcase{!A}{\ltrue}{$\indfrm$માં $0$ ડાબા અને
    $0$ જમણા કૌંસ છે.}}{}

\item \indcase{!A}{\Atom{R}{t_1,\dots,t_n}}{$l(\indfrm)=1+l(t_1)+
  \dots+l(t_n)=1+r(t_1)+\dots+r(t_n)=r(\indfrm)$. અહીં આપણે
  અભ્યાસપ્રશ્ન તરીકે છોડેલી હકીકત વાપરી છે કે દરેક પદ~$t$ માટે
  $l(t)=r(t)$.}

\item \indcase{!A}{\eq[t_1][t_2]}{$l(\indfrm)=l(t_1)+l(t_2)=
  r(t_1)+r(t_2)=r(\indfrm)$.}

\tagitem{prvNot}{\indcase{!A}{\lnot !B}{અનુમાન-પૂર્વધારણા મુજબ
    $l(!B)=r(!B)$. તેથી $l(\indfrm)=l(!B)=r(!B)=r(\indfrm)$.}}{}

\item \indcase{!A}{(!B \ast !C)}{અનુમાન-પૂર્વધારણા મુજબ $l(!B)=
  r(!B)$ અને $l(!C)=r(!C)$. તેથી $l(\indfrm)=1+l(!B)+l(!C)=
  1+r(!B)+r(!C)=r(\indfrm)$.}

\tagitem{prvAll}{\indcase{!A}{\lforall[x][!B]}{અનુમાન-પૂર્વધારણા
    મુજબ $l(!B)=r(!B)$. તેથી $l(\indfrm)=l(!B)=r(!B)=
    r(\indfrm)$.}}{}

% Print case for \lexists if prvEx and not prvAll
\tagitem{prvAll,notprvEx}{}{\indcase{!A}{\lexists[x][!B]}{અનુમાન-પૂર્વધારણા
    મુજબ $l(!B)=r(!B)$. તેથી $l(\indfrm)=l(!B)=r(!B)=
    r(\indfrm)$.}}

% Just say similarly if prvEx and prvAll
\tagitem{notprvAll,notprvEx}{}{\indcase{!A}{\lexists[x][!B]}{એ જ રીતે.}}
\end{enumerate}
\end{proof}

\begin{defn}[ઉચિત આરંભખંડ]
સંકેતોની શ્રેણી~$!B$ને સંકેતોની શ્રેણી~$!A$નો \emph{ઉચિત આરંભખંડ}
ત્યારે કહીએ જ્યારે $!B$ને કોઈ બિન-રિક્ત સંકેતશ્રેણી સાથે જોડવાથી
$!A$ મળે.
\end{defn}

\begin{lem}\ollabel{lem:no-prefix}
જો $!A$ !!a{formula} હોય અને $!B$, $!A$નો ઉચિત આરંભખંડ હોય, તો
$!B$ !!a{formula} નથી.
\end{lem}

\begin{proof}
અભ્યાસપ્રશ્ન.
\end{proof}

\begin{prob}
\olref[fol][syn][unq]{lem:no-prefix} સાબિત કરો.
\end{prob}

\begin{prop}
\ollabel{prop:unique-atomic}
જો $!A$ આણ્વિક !!{formula} હોય, તો તે નીચેનામાંથી એક અને માત્ર એક
શરત સંતોષે છે.
\begin{enumerate}
\tagitem{prvFalse}{$!A \ident \lfalse$.}{}
\tagitem{prvTrue}{$!A \ident \ltrue$.}{}
\item $!A \ident \Atom{R}{t_1,\dots,t_n}$, જ્યાં $R$ કોઈ $n$-સ્થાની
  !!{predicate}, $t_1$, \dots, $t_n$ પદો છે અને $R$, $t_1$, \dots,
  $t_n$માંથી દરેક એકમાત્ર રીતે નિર્ધારિત છે.
\item $!A \ident \eq[t_1][t_2]$, જ્યાં $t_1$ અને $t_2$ એકમાત્ર રીતે
  નિર્ધારિત પદો છે.
\end{enumerate}
\end{prop}

\begin{proof}
અભ્યાસપ્રશ્ન.
\end{proof}

\begin{prob}
\olref[fol][syn][unq]{prop:unique-atomic} સાબિત કરો. (સંકેત: પદો માટે
\olref[fol][syn][unq]{lem:no-prefix}નું રૂપ ઘડીને સાબિત કરો.)
\end{prob}

\begin{prop}[એકમાત્ર વાચનીયતા]
દરેક !!{formula} નીચેનામાંથી એક અને માત્ર એક શરત સંતોષે છે.
\begin{enumerate}
\item $!A$ આણ્વિક છે.

\tagitem{prvNot}{$!A$, $\lnot !B$ સ્વરૂપનું છે.}{}

\tagitem{prvAnd}{$!A$, $(!B \land !C)$ સ્વરૂપનું છે.}{}

\tagitem{prvOr}{$!A$, $(!B \lor !C)$ સ્વરૂપનું છે.}{}

\tagitem{prvIf}{$!A$, $(!B \lif !C)$ સ્વરૂપનું છે.}{}

\tagitem{prvIff}{$!A$, $(!B \liff !C)$ સ્વરૂપનું છે.}{}

\tagitem{prvAll}{$!A$, $\lforall[x][!B]$ સ્વરૂપનું છે.}{}

\tagitem{prvEx}{$!A$, $\lexists[x][!B]$ સ્વરૂપનું છે.}{}
\end{enumerate}
વળી, દરેક કિસ્સામાં $!B$, અથવા $!B$ અને $!C$, એકમાત્ર રીતે
નિર્ધારિત છે. તેનો અર્થ, દાખલા તરીકે, એવા જુદા યુગ્મો $!B$, $!C$ અને
$!B'$, $!C'$ નથી કે $!A$ બંને
\iftag{prvIf}{$(!B \lif !C)$ અને $(!B' \lif !C')$ સ્વરૂપનું હોય.}{
    \iftag{prvOr}{$(!B \lor !C)$ અને $(!B' \lor !C')$ સ્વરૂપનું હોય.}{
      \iftag{prvAnd}{$(!B \land !C)$ અને $(!B' \land !C')$ સ્વરૂપનું હોય.}{}}}
\end{prop}

\begin{proof}
રચના-નિયમો મુજબ કોઈ !!a{formula} આણ્વિક ન હોય તો તે આરંભકૌંસ~(,
\iftag{prvNot}{$\lnot$,}{} અથવા કોઈ પરિમાણકથી શરૂ થવું જોઈએ. બીજી
તરફ, નીચેના સંકેતોમાંથી કોઈથી શરૂ થતું દરેક !!{formula} આણ્વિક હોવું
જોઈએ: !!a{predicate}, !!a{function}, !!a{constant}
\iftag{prvFalse}{, $\lfalse$}{}\iftag{prvTrue}{, $\ltrue$}{}.

તેથી ખરેખર આપણે માત્ર એટલું બતાવવાનું છે કે જો $!A$, $(!B\ast !C)$
સ્વરૂપનું અને સાથે $(!B'\mathbin{\ast'}!C')$ સ્વરૂપનું હોય, તો
$!B\ident !B'$, $!C\ident !C'$ અને $\ast={\ast'}$.

ધારો કે $!A\ident(!B\ast !C)$ અને
$!A\ident(!B'\mathbin{\ast'}!C')$ બંને છે. હવે કાં તો
$!B\ident !B'$ અથવા નહિ. જો હોય તો સ્પષ્ટ છે કે $\ast={\ast'}$ અને
$!C\ident !C'$, કારણ કે ત્યાર પછી એ $!A$ની એક જ જગ્યાએ શરૂ થતી અને
સમાન લંબાઈ ધરાવતી ઉપશ્રેણીઓ છે. બીજો કિસ્સો $!B\not\ident !B'$ છે.
$!B$ અને $!B'$ બંને $!A$ની એક જ જગ્યાએ શરૂ થતી ઉપશ્રેણીઓ હોવાથી
એક બીજાનો ઉચિત આરંભખંડ હોવો જોઈએ. પરંતુ
\olref{lem:no-prefix} મુજબ એ અશક્ય છે.
\end{proof}


\end{document}
